\emph{(i)} For $n$ even every vector anticommutes with $\omega$, so every bivector, and hence every even element, commutes with it: $\omega$ is central in $\Cl^{0}(V)$. A basis blade $e_{B}$ commutes or anticommutes with each bivector $e_{i}e_{k}$, and distinct blades multiply a fixed blade to distinct blades, so the centre of $\Cl^{0}(V)$ is spanned by the even blades that commute with every bivector; these are $1$ alone for $n$ odd and $1, \omega$ for $n$ even (the check enumerates every signature with $n \le 8$). Since $\omega^{2} = e_{1}\cdots e_{n}e_{1}\cdots e_{n} = (-1)^{n(n-1)/2}\prod_{i}e_{i}^{2} = (-1)^{n(n-1)/2 + q}$ and $\rev{\omega} = (-1)^{n(n-1)/2}\omega$, for $n$ even and $\omega^{2} = 1$ the elements $P_{\pm} = \tfrac12(1 \pm \omega)$ are central idempotents with $P_{+}P_{-} = 0$, and $\rev{P}_{+} = P_{-}$ exactly when $n(n-1)/2$ is odd, that is $n \equiv 2 \pmod 4$. Then for $\psi_{1}, \psi_{2} \in \Cl^{0}(V)P_{+}$, using $\psi_{i} = \psi_{i}P_{+}$, $\rev{\psi}_{1} = P_{-}\rev{\psi}_{1}$ and the centrality of $P_{-}$ on the even element $\rev{\psi}_{1}\psi_{2}$,
\[
\rev{\psi}_{1}\psi_{2} \;=\; P_{-}\,\rev{\psi}_{1}\psi_{2}\,P_{+} \;=\; \rev{\psi}_{1}\psi_{2}\,P_{-}P_{+} \;=\; 0 .
\]
For a causal signature with $q = 1$ the residue is $n - 2$, and with $p = 1$ it is $2 - n$; either way $r \equiv 0 \pmod 4$ exactly when $n \equiv 2 \pmod 4$, and then the exponent $n(n-1)/2 + q$ is even (odd plus one for $q = 1$, odd plus $n - 1$ for $p = 1$), so $\omega^{2} = 1$. When $n \equiv 0 \pmod 4$ the same exponent is odd and $\omega^{2} = -1$, and the residue is $2$ or $6$ (Proposition \claimref{synthesis.oriented-causal-towers-mirror-pair-every-window}(i)). The check confirms the sign rules against the algebra for every signature with $n \le 8$ and applies them to the causal signatures for $n \le 64$. The fixed points of $r \mapsto -r$ modulo $8$ are $0$ and $4$. Residue, Morita class and division-ring type are constant along a tower (Fact \claimref{dark-matter.imported-eight-towers-mod-8-clock}); the statement here is about every causal member, not only the lowest.

\emph{(ii)} For a causal signature of residue other than $0$ and $4$, either $n$ is odd and $\Cl^{0}(V)$ is a semisimple real algebra with centre $\R$, hence simple, or $n \equiv 0 \pmod 4$ and the centre is $\R \oplus \R\omega$ with $\omega^{2} = -1$, a field, so again $\Cl^{0}(V)$ is simple. Write $A = \Cl^{0}(V)$. With $\psi_{1} = f$ and $\psi_{2} = af$, $a \in A$, the pairing is $\rev{f}\,a\,f$. In a simple ring $AfA = A$, so $\rev{f}AfA = \rev{f}A \ni \rev{f} \ne 0$ and therefore $\rev{f}Af \ne 0$: the pairing is not identically zero on $S$.

\emph{(iii)} With the bases of $S$ built from Definition \claimref{obj.dark-tower-matter.ideals-fields-momentum} exactly as in the check of Proposition \claimref{synthesis.lowest-causal-residues-zero-four-content} (dimensions $1, 2, 2, 4, 4, 8, 8, 8, 8$ at $(1,1)$, $(1,2)$, $(2,1)$, $(1,3)$, $(3,1)$, $(1,4)$, $(4,1)$, $(1,5)$, $(5,1)$), the check evaluates $\rev{s}_{i}s_{l}$ on every pair of basis elements; it is zero for all pairs exactly at $(1,1)$ (either ideal, $f$ or $1 - f$, which are $P_{\pm}$ there) and at $(1,5)$, $(5,1)$ for either chiral factor, and nonzero for some pair at the six other signatures. By bilinearity, vanishing on basis pairs is vanishing on $S \times S$.

\emph{(iv)} On every ordered pair of basis elements of $S$ at $(1,2)$, $(2,1)$, $(1,3)$, $(3,1)$ the check confirms $m^{2}(j(s_{i}) + j(s_{l})) = 2\ip{xx^{\dagger}}_{0}$ with $x = \rev{s}_{i}s_{l}$ and $x^{\dagger} = e_{t}\rev{x}e_{t}^{-1}$ (Theorem \claimref{dark-tower-matter.where-mass-begins}(i), whose proof is a polynomial identity), and counts the pairs with nonzero mass ($2$, $2$, $8$, $8$). At $(1,5)$ and $(5,1)$, by (i), $x = 0$ for every pair, so the right-hand side is zero, while the momentum sum of the first two basis elements has $m^{2} = \tfrac14$, the value of Proposition \claimref{synthesis.lowest-causal-residues-zero-four-content}(iii). Hence the identity fails in dimension six. In dimension five Theorem \claimref{dark-tower-matter.where-mass-begins}(ii) gives the one-constituent mass as $\tfrac12\ip{xx^{\dagger}}_{0}$ with $x = \rev{\psi}\psi$, so in each of the populated causal towers (dimensions three, four, five) the mass the store has proved is a function of the reversion pairing. The momentum reading at $(1,5)$ rests on Premise H1 extended to dimension six, as recorded in Proposition \claimref{synthesis.lowest-causal-residues-zero-four-content}; part (iv) inherits that extension and parts (i) to (iii) do not use it.

\emph{The consequence.} With one spinor ideal per tower, a reading in which a composite is massive only through a nonvanishing pairing of its constituents gives no massive matter at any causal member of residues $0$ and $4$, in any dimension, by (i), and leaves the six other residues their pairing, by (ii); the momentum reading gives mass at $(1,5)$ and $(5,1)$ (Proposition \claimref{synthesis.lowest-causal-residues-zero-four-content}(iii)). The flag bit exists exactly at $n \equiv 0 \pmod 4$ (Proposition \claimref{synthesis.oriented-causal-towers-mirror-pair-every-window}), so residues $0$, $4$ and every odd residue lack it; since the odd residues have a nonvanishing pairing by (ii) and are taken as populated, absence of the bit cannot be the reason for the emptiness of $0$ and $4$. The pairing reading is a costed clause, chosen with the fit known (Proposition \claimref{dark-matter.status-rule}). Filed alongside, synthesis.residues-zero-four-real-chirality-anomaly cites (i) as its lemma; under its hypothesis (G6), that the residue-$4$ tower's own gauge fields are fields on its six-dimensional space acting on its constituents, the anomaly of a single chiral ideal excludes residue-$4$ matter on either reading, which closes the momentum reading's escape there.

(Reading, parenthetical and from memory: in the field-theory language this is the statement that in dimensions $2 \bmod 4$ the Lorentz-invariant spinor bilinears pair opposite chiralities only, charge conjugation preserving chirality there, so a single-chirality field admits no mass term; the two-constituent momentum mass is then the invariant mass of a two-particle state of massless quanta, which any non-collinear pair has. The statement above does not rely on this reading.)
